\chapter{De kwantummechanische harmonische oscillator}\label{ch:b3:harmonic-oscillator}

Vrijwel niets in de natuur is precies een harmonische oscillator, en
vrijwel alles is er bij benadering een: elk stelsel dat uit een stabiel
evenwicht wordt geduwd --- de binding van een molecuul, een atoom in een
kristal, een brug, een modus van het elektromagnetische veld --- voelt een
terugdrijvende kracht evenredig met de uitwijking, want elke gladde
potentiaal is een parabool op de bodem van haar put. Wie de
kwantummechanische oscillator één keer oplost, lost dus de \omterm{prop:b3:lagrangian-mechanics:modes}{kleine
trillingen} van de hele wereld op. Dit hoofdstuk lost hem op in de
algebraïsche stijl die het handschrift van de kwantummechanica is geworden:
twee \omterm{def:b3:harmonic-oscillator:ladder}{ladderoperatoren} klimmen en dalen langs een volmaakt gelijkmatige
trap van niveaus $\big(n +
\tfrac12\big)\hbar\omega$, waarbij de halve
trede onderaan --- de \omterm{thm:b3:harmonic-oscillator:spectrum}{nulpuntsenergie} --- een stelling is en geen keuze.
De gevolgen reiken van de vraag waarom helium nooit bevriest tot de manier
waarop een molecuul koolstofdioxide, dat op zijn eigen $\hbar\omega$
rinkelt, de uitgaande warmte van de aarde onderschept.

\section{De ladder}

\begin{definition}[Ladderoperatoren]\label{def:b3:harmonic-oscillator:ladder}
\index{ladderoperatoren}\index{annihilatieoperator}\index{creatieoperator}
Voer voor $\hat H = \hat p^2/2m + \tfrac12 m\omega^2\hat x^2$ het
dimensieloze, niet-hermitische paar
\[
\hat a = \sqrt{\frac{m\omega}{2\hbar}}\Big(\hat x +
\frac{\iu\hat p}{m\omega}\Big) , \qquad
\hat a^\dagger = \sqrt{\frac{m\omega}{2\hbar}}\Big(\hat x -
\frac{\iu\hat p}{m\omega}\Big) .
\]
in. Uit $[\hat x, \hat p] = \iu\hbar$ volgt
\[
[\hat a, \hat a^\dagger] = 1 , \qquad
\hat H = \hbar\omega\Big(\hat N + \tfrac12\Big) , \quad
\hat N = \hat a^\dagger\hat a .
\]
$\hat N$ is hermitisch; haar \omterm{def:b3:quantum-formalism:observable}{eigenwaarden} zullen \emph{quanta} tellen, en
$\hat a$ en $\hat a^\dagger$ --- de \emph{annihilatie}- en de
\emph{creatie}operator --- zullen er telkens één weghalen en toevoegen.
\end{definition}

\begin{theorem}[Het spectrum, met algebra alleen]\label{thm:b3:harmonic-oscillator:spectrum}
\index{nulpuntsenergie}
De \omterm{def:b3:quantum-formalism:observable}{eigenwaarden} van $\hat H$ zijn precies
\[
E_n = \Big(n + \tfrac12\Big)\hbar\omega , \qquad n = 0, 1, 2, \dots
\]
--- een oneindige ladder met gelijke treden $\hbar\omega$ boven een
grondniveau dat \emph{niet} nul is: de \emph{\omterm{thm:b3:harmonic-oscillator:spectrum}{nulpuntsenergie}}
$\tfrac12\hbar\omega$. De genormeerde eigentoestanden $\ket n$ hangen
samen via
\[
\hat a\ket n = \sqrt n\,\ket{n - 1} , \qquad
\hat a^\dagger\ket n = \sqrt{n + 1}\,\ket{n + 1} ,
\qquad
\ket n = \frac{(\hat a^\dagger)^n}{\sqrt{n!}}\,\ket0 .
\]
\end{theorem}

\begin{proof}
Uit de \omterm{def:b3:quantum-formalism:commutator}{commutator} volgt $\hat N\hat a = \hat a(\hat N - 1)$: heeft
$\ket\nu$ de $\hat N$-eigenwaarde $\nu$, dan is $\hat a\ket\nu$ een
eigenvector met $\nu - 1$ (of de nulvector). Elke afdaling is alleen
toegestaan zolang $\nu \ge 0$, want $\nu = \braket{\nu}{\hat
N\nu} = \|\hat a\ket\nu\|^2 \ge 0$; de afdaling moet dus eindigen, en zij
eindigt alleen op een toestand met $\hat a\ket{\nu_0}
= 0$, waaruit
$\nu_0 = 0$. Dus doorloopt $\nu$ de niet-negatieve gehele getallen. De
normeringen volgen uit $\|\hat a\ket n\|^2 = n$ en $\|\hat a^\dagger\ket n\|^2 = n + 1$.
\end{proof}

\begin{omfigure}
\begin{tikzpicture}[font=\small, >=Stealth]
  \draw[->] (0,-0.3) -- (0,5.4);
  \node[anchor=south] at (0,5.4) {$E$};
  \foreach \n in {0,1,2,3,4}
    {\draw[very thick, omDef] (0.5,{0.55+\n*1.05}) -- (3.3,{0.55+\n*1.05});
     \node[anchor=west] at (3.45,{0.55+\n*1.05})
       {$\ket{\n}$\; $E = (\n + \tfrac12)\hbar\omega$};}
  \draw[->, omThm, very thick] (1.35,0.62) -- (1.35,1.53);
  \node[omThm, anchor=east] at (1.28,1.08) {$\hat a^\dagger$};
  \draw[->, omExo, very thick] (2.45,1.53) -- (2.45,0.62);
  \node[omExo, anchor=west] at (2.52,1.08) {$\hat a$};
  \draw[<->] (0.15,0) -- (0.15,0.55);
  \node[anchor=west, font=\footnotesize] at (0.28,0.22) {$\tfrac12\hbar\omega$};
  \draw[dashed] (0,0) -- (3.3,0);
\end{tikzpicture}

{\small De ladder van de oscillator: gelijke treden $\hbar\omega$,
beklommen door $\hat a^\dagger$ en afgedaald door $\hat a$, staand op een
vloer die een halve trede boven de klassieke \omterm{def:b3:relativistic-dynamics:momentum-energy}{rustenergie} ligt.}
\end{omfigure}

\begin{proposition}[De toestanden in de ruimte]\label{prop:b3:harmonic-oscillator:wavefunctions}
De grondtoestand is de gaussfunctie
\[
\varphi_0(x) = \Big(\frac{m\omega}{\pi\hbar}\Big)^{1/4}
\eu^{-m\omega x^2/2\hbar} ,
\]
die men vindt door de \emph{eerste-orde}vergelijking $\hat a\varphi_0
= 0$
op te lossen; zij verzadigt de ongelijkheid van Heisenberg, $\Delta x\,\Delta p =
\hbar/2$, met $\Delta x = \sqrt{\hbar/2m\omega}$. Herhaald
toepassen van $\hat a^\dagger$ brengt $\varphi_n$ voort: een veelterm van
graad $n$ (met $n$ knopen) maal dezelfde gaussfunctie. Voor grote $n$
oscilleert $|\varphi_n|^2$ snel rond de klassieke verblijftijdverdeling,
die het grootst is nabij de keerpunten --- het correspondentiebeginsel in
één plaatje.
\end{proposition}

\begin{proof}[Gedeeltelijk bewijs]
$\hat a\varphi_0 = 0$ luidt $\varphi_0' = -(m\omega/\hbar)x
\varphi_0$: de gaussfunctie, genormeerd. Verzadiging: de gaussfunctie is
het gelijkheidsgeval van \cref{thm:b3:quantum-formalism:uncertainty}. De
veeltermstructuur volgt eruit dat $\hat a^\dagger$ van de eerste orde is
in $x$ en $\dd/\dd x$; de uitspraak voor grote $n$ wordt nagegaan in
\cref{exo:b3:harmonic-oscillator:12}.
\end{proof}

\begin{omfigure}
\begin{tikzpicture}
  \begin{axis}[omaxis, axis x line=bottom, axis y line=left,
    width=6.6cm, height=5.2cm, xmin=-4, xmax=4, ymin=0, ymax=3.6,
    xlabel={$x$ (in eenheden $\sqrt{\hbar/m\omega}$)}, ylabel={},
    xlabel style={at={(axis description cs:0.5,-0.1)}, anchor=north},
    xtick={-2,0,2}, ytick=\empty, clip=false]
    \addplot[omDef, thick, domain=-4:4, samples=90]
      {0.56*exp(-x*x)+0.25};
    \addplot[omThm, thick, domain=-4:4, samples=90]
      {1.13*x*x*exp(-x*x)+1.3};
    \addplot[omExo, thick, domain=-4:4, samples=120]
      {0.28*(2*x*x-1)^2*exp(-x*x)+2.35};
    \node[omDef, font=\scriptsize, anchor=west] at (axis cs:2.4,0.5) {$|\varphi_0|^2$};
    \node[omThm, font=\scriptsize, anchor=west] at (axis cs:2.4,1.55) {$|\varphi_1|^2$};
    \node[omExo, font=\scriptsize, anchor=west] at (axis cs:2.4,2.6) {$|\varphi_2|^2$};
  \end{axis}
  \begin{axis}[omaxis, axis x line=bottom, axis y line=left,
    width=6.6cm, height=5.2cm, xmin=-5.2, xmax=5.2, ymin=0, ymax=0.45,
    xshift=7.1cm,
    xlabel={$x$}, ylabel={},
    xlabel style={at={(axis description cs:0.5,-0.1)}, anchor=north},
    xtick=\empty, ytick=\empty, clip=false]
    \addplot[omDef, domain=-4.6:4.6, samples=300]
      {(64*x^6-480*x^4+720*x^2-120)^2*exp(-x*x)/81680};
    \addplot[omThm, very thick, dashed, domain=-3.54:3.54, samples=120]
      {0.3183/sqrt(max(0.35,13-x*x))};
    \node[omThm, font=\scriptsize, anchor=south] at (axis cs:0,0.36) {klassieke verblijftijd};
    \node[omDef, font=\scriptsize, anchor=west] at (axis cs:1.9,0.42) {$|\varphi_6|^2$};
  \end{axis}
\end{tikzpicture}

{\small Links: de laagste kansdichtheden (verticaal verschoven): een
knooploze gaussfunctie, daarna $n$ knopen voor $\varphi_n$. Rechts:
bij grote $n$ oscilleert de kwantumdichtheid rond de klassieke verdeling,
die zich ophoopt bij de keerpunten waar de trillende massa talmt.}
\end{omfigure}

\section{De nulpuntsenergie is echt}

\begin{remark}[Waarom de vloer niet lager kan]\label{rem:b3:harmonic-oscillator:zeropoint}
Een toestand met energie nul zou $\langle\hat p^2\rangle =
\langle\hat x^2\rangle = 0$ vergen: volmaakt stil \emph{en} volmaakt gecentreerd, wat
$\Delta x\,\Delta p \ge \hbar/2$ verbiedt. De grondtoestand is het beste
compromis dat de ongelijkheid toelaat, en $\tfrac12\hbar\omega$ is de
huur. Het is geen boekhoudconstante: de nulpuntsbeweging vervaagt bij het
absolute nulpunt de patronen van röntgendiffractie; zij geeft lichtere
isotopen zwakkere effectieve bindingen (H$_2$ en D$_2$ dissociëren bij
meetbaar verschillende energieën uit dezelfde elektronische put); en in
helium is zij zo heftig --- lichte atomen, zwakke aantrekking --- dat de
vloeistof onder haar eigen dampdruk \emph{nooit} bevriest: het enige
element dat bij het absolute nulpunt nog vloeibaar is en pas onder 25
atmosfeer hulp stolt.
\end{remark}

\begin{example}[Schalen van $\hbar\omega$]\label{ex:b3:harmonic-oscillator:scales}
De binding van koolstofmonoxide ($\omega = \qty{4.1e14}{rad/s}$):
$\hbar\omega =
\qty{0.27}{eV}$, tien keer de $k_{\text{B}}T$ van
kamertemperatuur --- moleculaire trillingen liggen in alledaagse lucht
bevroren, en daarom baarden de soortelijke warmten van tweeatomige gassen
de negentiende eeuw zorgen. Een slinger ($\omega = \qty{5}{rad/s}$):
$\hbar\omega = \qty{3e-15}{eV}$, hopeloos buiten elk oplossend vermogen
--- de correspondentielimiet. Een spiegel van LIGO van \qty{40}{kg},
zo opgehangen dat zij bij \qty{100}{Hz} zwaait, heeft als oscillatormodus de nulpuntsamplitude $\Delta x =
\sqrt{\hbar/2m\omega} \approx \qty{5e-20}{m}$ --- en het observatorium
lost routinematig verplaatsingen \emph{op} deze kwantumvloer op: de
nulpuntsbeweging van een voorwerp van veertig kilogram is inmiddels een
technische randvoorwaarde (\cref{exo:b3:harmonic-oscillator:3}).
\end{example}

\section{Coherente toestanden: het klassieke gezicht van de kwantumoscillator}

\begin{definition}[Coherente toestanden]\label{def:b3:harmonic-oscillator:coherent}
\index{coherente toestand}
Een \emph{coherente toestand} $\ket\alpha$ is een eigenvector van de
\omterm{def:b3:harmonic-oscillator:ladder}{annihilatieoperator}, $\hat a\ket\alpha = \alpha\ket\alpha$, met $\alpha$
een willekeurig complex getal. Ontwikkeld op de ladder:
\[
\ket\alpha = \eu^{-|\alpha|^2/2}\sum_{n=0}^\infty
\frac{\alpha^n}{\sqrt{n!}}\,\ket n :
\]
het aantal quanta is poissonverdeeld met gemiddelde $\langle\hat
N\rangle = |\alpha|^2$ en spreiding $\Delta N = |\alpha|$.
\end{definition}

\begin{proposition}[Waarom lasers klassiek zijn]\label{prop:b3:harmonic-oscillator:coherent-motion}
Een \omterm{def:b3:harmonic-oscillator:coherent}{coherente toestand} is de gaussfunctie van de grondtoestand, verschoven
in de \omterm{def:b3:hamiltonian-mechanics:hamiltonian}{faseruimte}; onder de evolutie van de oscillator \emph{blijft} zij
coherent, met $\alpha(t) = \alpha\,\eu^{-\iu\omega t}$: haar middelpunt
voert precies de klassieke beweging uit, $\langle\hat x\rangle(t) =
x_0\cos\omega t + \cdots$, terwijl haar breedten voor altijd de minimale
$\Delta x\,\Delta p = \hbar/2$ blijven --- een golfpakket dat nooit
uitvloeit. \omterm{def:b3:harmonic-oscillator:coherent}{Coherente toestanden} zijn hoe een kwantumoscillator een
klassieke nabootst; het licht van een laser is een \omterm{def:b3:harmonic-oscillator:coherent}{coherente toestand} van
een veldmodus, met een poissonverdeeld aantal fotonen (de \emph{hagelruis}
van elke fotodetector) en een veld dat trilt zoals Maxwell zei.
\end{proposition}

\begin{proof}[Gedeeltelijk bewijs]
De ontwikkeling volgt door $\ket\alpha = \sum c_n\ket n$ in $\hat a\ket\alpha = \alpha\ket\alpha$ te schrijven: $c_n = \alpha c_{n-1}/\sqrt
n$. Evolutie: elke $\ket n$ pikt $\eu^{-\iu(n + 1/2)\omega t}$ op, wat
weer optelt tot een \omterm{def:b3:harmonic-oscillator:coherent}{coherente toestand} van $\alpha\eu^{-\iu\omega t}$ (op
een globale fase na). $\langle\hat x\rangle \propto
\operatorname{Re}\alpha(t)$ volgt uit $\hat x \propto \hat a +
\hat a^\dagger$. Dat de toestand de verschoven gaussfunctie is, wordt hier
aangenomen.
\end{proof}

\begin{omfigure}
\begin{tikzpicture}[font=\small, >=Stealth]
  \draw[->] (-3.1,0) -- (3.4,0) node[right] {$\langle\hat x\rangle$};
  \draw[->] (0,-3.0) -- (0,3.2) node[above] {$\langle\hat p\rangle$};
  \draw[gray, dashed] (0,0) circle (2.2);
  \fill[omDef!25, draw=omDef, thick] (2.2,0) circle (0.42);
  \fill[omDef!25, draw=omDef, thick] (50:2.2) circle (0.42);
  \fill[omDef!25, draw=omDef, thick] (100:2.2) circle (0.42);
  \draw[->, thick] (2.53,0.9) arc (20:40:2.6);
  \node[anchor=west] at (2.62,1.35) {$\omega$};
  \fill[omThm!30, draw=omThm, thick] (0,0) circle (0.42);
  \node[omThm, font=\footnotesize, anchor=east] at (-0.3,-0.68) {$\ket0$};
  \node[omDef, font=\footnotesize, anchor=west] at (2.3,-0.62) {$\ket\alpha$};
  \node[font=\footnotesize, anchor=west] at (-3.05,2.75)
    {minimale vlek $\Delta x\,\Delta p = \hbar/2$, klassiek rondgaand};
\end{tikzpicture}

{\small Portret in de \omterm{def:b3:hamiltonian-mechanics:hamiltonian}{faseruimte} (vergelijk
\cref{ch:b3:hamiltonian-mechanics}): de grondtoestand is een vlek met
minimale onzekerheid in de oorsprong; een \omterm{def:b3:harmonic-oscillator:coherent}{coherente toestand} is dezelfde
vlek, verschoven, die met $\omega$ rondgaat zonder te vervormen --- de
beste nabootsing van een klassieke trilling die de kwantummechanica kent.}
\end{omfigure}

\begin{method}[Algebra van de oscillator]\label{met:b3:harmonic-oscillator:method}
(1) Druk uit wat er gevraagd wordt in $\hat a$ en $\hat a^\dagger$: $\hat x
= \sqrt{\hbar/2m\omega}\,(\hat a + \hat a^\dagger)$, $\hat p =
\iu\sqrt{m\hbar\omega/2}\,(\hat a^\dagger - \hat a)$. (2) Schuif de $\hat a$'s met $[\hat a, \hat a^\dagger] = 1$ naar rechts; gebruik $\hat a\ket0 = 0$. (3) Matrixelementen: $\hat x$ verbindt alleen naburige treden --- de
selectieregel $\Delta n = \pm1$ van de trillingsspectroscopie. (4) Elke
kwadratische \omterm{def:b3:hamiltonian-mechanics:hamiltonian}{hamiltoniaan} (gekoppelde oscillatoren, kringen, veldmodi)
valt uiteen in onafhankelijke ladders: bepaal eerst de \omterm{prop:b3:lagrangian-mechanics:modes}{eigentrillingen} en
kwantiseer elk ervan. (5) Eerst getallen: $\hbar\omega$ tegen
$k_{\text{B}}T$ beslist of het stelsel kwantummechanisch of klassiek is,
nog voordat er enige algebra aan te pas komt.
\end{method}

\section{Opgaven}

\begin{exercise}[$\star$]\label{exo:b3:harmonic-oscillator:1}
(a) Ga $[\hat a, \hat a^\dagger] = 1$ na uit $[\hat x, \hat p] =
\iu\hbar$. (b) Ga $\hat H = \hbar\omega(\hat a^\dagger\hat a +
\tfrac12)$ na door rechtstreeks uit te werken. (c) Keer om en druk $\hat x$ en $\hat p$ uit. (d) Toon aan dat $\hat N$ hermitisch is en leg uit
waarom $\hat a$ op zichzelf geen observabele kan zijn.
\end{exercise}

\begin{exercise}[$\star$]\label{exo:b3:harmonic-oscillator:2}
Op de eigentoestand $\ket n$: (a) toon aan dat $\langle\hat x\rangle =
\langle\hat p\rangle = 0$; (b) bereken $\langle\hat x^2\rangle$ en
$\langle\hat p^2\rangle$; (c) leid $\Delta x\,\Delta p = (n +
\tfrac12)\hbar$ af; (d) ga de viriaalverhouding $\langle E_k\rangle =
\langle E_p\rangle$ na.
\end{exercise}

\begin{exercise}[$\star$]\label{exo:b3:harmonic-oscillator:3}
Nulpuntsamplitudes $\Delta x = \sqrt{\hbar/2m\omega}$: bereken die voor
(a) het CO-molecuul ($\mu = \qty{1.14e-26}{kg}$, $\omega =
\qty{4.1e14}{rad/s}$), vergeleken met de bindingslengte \qty{0.11}{nm};
(b) een waterstofatoom in een vaste stof ($m =
\qty{1.7e-27}{kg}$,
$\hbar\omega = \qty{0.1}{eV}$); (c) een spiegel van LIGO ($m = \qty{40}{kg}$, $f = \qty{100}{Hz}$); (d) rangschik de drie als fracties
van de afmetingen van hun stelsels en geef commentaar op wie
``kwantummechanisch'' is.
\end{exercise}

\begin{exercise}[$\star$]\label{exo:b3:harmonic-oscillator:4}
(a) Bereken met de ladderbetrekkingen $\bra m\hat x\ket n$ en toon aan dat
het nul is tenzij $m = n \pm 1$. (b) Leid de trillingsselectieregel
$\Delta n = \pm1$ voor lichtabsorptie af (de koppeling is
$\propto\hat x$). (c) Waarom absorbeert een heteronucleair molecuul (CO)
infrarood licht en N$_2$ niet? (d) Echte moleculen vertonen zwakke
``boventoon''-lijnen bij $\approx 2\hbar\omega$: wat onthult dat over de
potentiaal?
\end{exercise}

\begin{exercise}[$\star\star$]\label{exo:b3:harmonic-oscillator:5}
(a) Los $\hat a\varphi_0 = 0$ op als differentiaalvergelijking en normeer.
(b) Breng $\varphi_1$ en $\varphi_2$ voort door $\hat a^\dagger$ toe te
passen. (c) Ga $\varphi_1 \perp \varphi_0$ na op grond van de pariteit
alleen. (d) Schets de drie dichtheden en tel de knopen na.
\end{exercise}

\begin{exercise}[$\star\star$]\label{exo:b3:harmonic-oscillator:6}
Vooruitblik op Boltzmann. Een verzameling identieke oscillatoren bij
temperatuur $T$ bezet niveau $n$ met kans $\propto
\eu^{-E_n/k_{\text{B}}T}$ (volume van jaar 2, boltzmannfactor). (a) Toon
aan dat het gemiddelde kwantumgetal $\langle n\rangle =
1/(\eu^{\hbar\omega/k_{\text{B}}T} - 1)$ is. (b) Reken dit uit voor de
CO-trilling bij \qty{300}{K} en bij \qty{2000}{K}. (c) Toon aan dat de
gemiddelde energie bij hoge temperatuur naar $k_{\text{B}}T$ gaat
(equipartitie teruggevonden) en bij lage naar $\tfrac12\hbar\omega$. (d)
Bij welke temperatuur bevat een trilling van \qty{15}{\micro m} (de buiging
van koolstofdioxide) $\langle n\rangle = 0.1$?
\end{exercise}

\begin{exercise}[$\star\star$]\label{exo:b3:harmonic-oscillator:7}
Statistiek van \omterm{def:b3:harmonic-oscillator:coherent}{coherente toestanden}. (a) Toon met de ontwikkeling van
$\ket\alpha$ aan dat $\mathcal P(n)$ poissonverdeeld is met gemiddelde
$|\alpha|^2$. (b) Toon aan dat $\langle\hat N\rangle = |\alpha|^2$ en
$\Delta N = |\alpha|$. (c) Een laser van \qty{1}{mW} bij \qty{633}{nm}:
fotonen per seconde, en de relatieve fluctuatie $\Delta N/\langle N\rangle$ in één seconde. (d) Hagelruis: toon aan dat de verhouding van
ruis tot signaal van de fotostroom als $1/\sqrt{\langle N\rangle}$ daalt
--- waarom heldere bundels er glad uitzien.
\end{exercise}

\begin{exercise}[$\star\star$]\label{exo:b3:harmonic-oscillator:8}
Evolutie van een \omterm{def:b3:harmonic-oscillator:coherent}{coherente toestand}. (a) Pas de evolutiefasen op de
ontwikkeling toe en toon $\ket{\alpha(t)}$ aan met $\alpha(t) =
\alpha\eu^{-\iu\omega t}$ (op een globale fase na). (b) Leid
$\langle\hat x\rangle(t)$ en $\langle\hat p\rangle(t)$ af en vergelijk met
de klassieke oplossing. (c) Waarom beschrijft een energie-eigentoestand,
hoewel stationair, \emph{geen} zwaaiende slinger --- welk kenmerk van de
\omterm{def:b3:harmonic-oscillator:coherent}{coherente toestand} doet dat wel? (d) Schat $|\alpha|$ voor een echte
slinger (\qty{10}{g}, amplitude \qty{10}{cm}, \qty{1}{Hz}) en geef
commentaar.
\end{exercise}

\begin{exercise}[$\star\star$]\label{exo:b3:harmonic-oscillator:9}
De kwantummechanische LC-kring. Een supergeleidende lus met $L =
\qty{10}{nH}$ en $C = \qty{0.4}{pF}$ laat lading en flux oscilleren zoals
$x$ en $p$. (a) Haar resonantiefrequentie (volume van jaar 1) en het
kwantum $\hbar\omega$ in $\unit{\micro eV}$. (b) Onder welke temperatuur
is $k_{\text{B}}T \ll \hbar\omega$, zodat de kring in haar grondtoestand
zit? (c) Waarom worden supergeleidende qubits in verdunningskoelkasten bij
\qty{20}{mK} bedreven? (d) De nulpuntsfluctuatie van de spanning $\Delta V = \Delta q/C$ met $\Delta q =
\sqrt{\hbar\omega C/2}$: reken haar uit.
\end{exercise}

\begin{exercise}[$\star\star\star$]\label{exo:b3:harmonic-oscillator:10}
Van der Waals uit de nulpuntsbeweging. Modelleer twee neutrale atomen op
afstand $R$ als twee identieke dipooloscillatoren (lading $e$, massa $m$,
frequentie $\omega_0$) waarvan de uitwijkingen koppelen via de
dipoolenergie $\lambda\,x_1x_2$ met $\lambda = e^2/2\pi\varepsilon_0R^3$.
(a) Toon aan dat de eigencoördinaten $(x_1 \pm x_2)/\sqrt2$ trillen bij
$\omega_\pm = \omega_0\sqrt{1 \pm \lambda/m\omega_0^2}$. (b) De
grondenergie is $\tfrac\hbar2(\omega_+ + \omega_-)$: ontwikkel tot tweede
orde in $\lambda$ en toon aan dat de wisselwerkingsenergie
\[
\Delta E = -\frac{\hbar\lambda^2}{8m^2\omega_0^3}
\ \propto\ -\frac{1}{R^6} .
\]
bedraagt. (c) Waarom is deze aantrekking universeel --- aanwezig zelfs
tussen atomen zonder enig permanent dipoolmoment? (d) De wet $1/R^6$ is de
vanderwaalskracht uit de correcties op reële gassen van het volume van
jaar 1: wat ``fluctueert'' er microscopisch --- en wat zou er met deze
kracht gebeuren in een wereld met $\hbar = 0$?
\end{exercise}

\begin{exercise}[$\star\star\star$]\label{exo:b3:harmonic-oscillator:11}
Zijbanden van een gevangen ion. Een ion in een harmonische val ($f =
\qty{1}{MHz}$) absorbeert laserlicht. Omdat het ion beweegt, vertoont zijn
absorptiespectrum de elektronische lijn bij $\nu_0$ geflankeerd door
lijnen bij $\nu_0 \pm f$: overgangen die het kwantumgetal van de
\emph{beweging} met $\mp1$ veranderen naast het elektronische. (a) Hoe
groot is $\hbar\omega_{\text{val}}$ in $\unit{neV}$, en waarom vergt het
oplossen van zijbanden een zeer smalle lijn? (b) De lijn $\nu_0 - f$
aandrijven haalt per cyclus één bewegingskwantum weg: leg
\emph{zijbandkoeling} tot in de grondtoestand uit. (c) Zodra $\langle n\rangle \approx
0$ verdwijnt de onderste zijband volledig --- waarom is
die asymmetrie een bewijs dat de kwantumgrondtoestand is bereikt? (d) Zo
wordt de bewegingsgrondtoestand van één enkel atoom --- en, langs andere
weg, van spiegels van LIGO op kilogramschaal --- gecertificeerd: zeg welke
``temperatuur'' het ion heeft bereikt voor $f = \qty{1}{MHz}$ en $\langle n\rangle = 0.05$.
\end{exercise}

\begin{exercise}[$\star\star\star$]\label{exo:b3:harmonic-oscillator:12}
De klassieke limiet, kwantitatief. Een klassieke oscillator met amplitude
$A$ brengt in $[x, x + \dd x]$ de fractie $\dd
t/T = \dd x/\pi\sqrt{A^2 - x^2}$ door. (a) Leid deze verblijftijdverdeling af. (b) Neem voor de
kwantumtoestand $n$ de $A_n$ uit $E_n =
\tfrac12 m\omega^2A_n^2$ en
vergelijk de klassieke verdeling met de (gegeven) plaatselijk uitgemiddelde
$|\varphi_n|^2$: waar stemmen zij overeen en waar moeten zij verschillen?
(c) Toon aan dat de relatieve afstand tussen naburige niveaus, $\Delta E/E$, als $1/n$ naar nul gaat: de energie wordt in de praktijk continu.
(d) Schat voor de slinger van \cref{exo:b3:harmonic-oscillator:8}(d) de
waarden van $n$ en $\Delta
E/E$, en rond het correspondentieargument in
één zin af.
\end{exercise}

\section{Vraagstuk: Het molecuul dat de aarde verwarmt}

\begin{problem}[{Weekendvraagstuk --- de kwantumladder van
koolstofdioxide en het broeikaseffect}]\label{pb:b3:harmonic-oscillator:1}
Een molecuul koolstofdioxide is een lineaire keten O=C=O: enkele
gekwantiseerde oscillatoren. Dat de $\hbar\omega$ van zijn buigmodus
toevallig midden in de uitgaande warmtegloed van de aarde ligt, is de
reden dat dit spoorgas --- vier moleculen op de tienduizend --- het
klimaat van de planeet stuurt. Gegevens: golfgetal van de buiging
$\tilde\nu_2 = \qty{667}{cm^{-1}}$ (de eenheid van de spectroscopist: $E = hc\tilde\nu$, $\qty{1}{cm^{-1}} =
\qty{1.24e-4}{eV}$); asymmetrische rek
$\tilde\nu_3 =
\qty{2349}{cm^{-1}}$; symmetrische rek $\tilde\nu_1 =
\qty{1388}{cm^{-1}}$; $k_{\text{B}}T$ bij $\qty{288}{K}$ is
\qty{24.8}{meV}; de wet van Wien (volume van jaar 2): $\lambda_{\max}T =
\qty{2898}{\micro m.K}$.

\medskip\noindent\textbf{Deel I --- De modi van een lineair molecuul.}
\begin{enumerate}
  \item Drie atomen hebben negen \omterm{def:b3:lagrangian-mechanics:coordinates}{vrijheidsgraden}: hoeveel daarvan zijn
        translaties van het hele molecuul, hoeveel rotaties (het molecuul
        is \emph{lineair}), en hoeveel trillingen blijven er over?
  \item Beschrijf de vier trillingen: symmetrische rek, asymmetrische rek
        en een tweevoudig ontaarde buiging --- waarom komt de buiging
        tweemaal voor?
  \item Licht koppelt aan een trillend \emph{elektrisch dipoolmoment}
        (\cref{exo:b3:harmonic-oscillator:4}). Welke modi van O=C=O
        moduleren het dipoolmoment, en welke is stil in het infrarood?
  \item Zet de drie golfgetallen om in fotongolflengten en plaats ze:
        welke liggen in het thermische infrarood?
  \item Waarom absorberen de twee belangrijkste luchtgassen, N$_2$ en
        O$_2$, vrijwel geen infrarood --- en met welk gevolg voor de
        doorzichtigheid van de atmosfeer?
  \item Waterdamp, gebogen en dipolair, absorbeert over een groot deel van
        het infrarood: in welk spectraal ``venster'' werkt de buiging van
        koolstofdioxide grotendeels alleen (vergelijk \qty{15}{\micro m}
        met de sterke waterbanden onder \qty{8}{\micro m} en boven
        \qty{20}{\micro m})?
\end{enumerate}

\medskip\noindent\textbf{Deel II --- De kwantumladder van de buiging.}
\begin{enumerate}[resume]
  \item Bereken $\hbar\omega_2$ in meV, en de ladder $E_n$.
  \item Welk deel van de moleculen bezet $n = 1$ bij \qty{288}{K} (ten
        opzichte van $n = 0$, boltzmannfactor)? En $n = 2$?
  \item Welke fotongolflengte drijft $n = 0 \to 1$ aan? Waarom
        overheerst \emph{dezelfde} golflengte bij de emissie?
  \item Rechtvaardig uit de harmonische ladder dat één golflengte de hele
        ladder bedient ($n \to n + 1$ voor elke $n$): welke eigenschap van
        de niveauafstand is aan het werk?
  \item Het molecuul draait ook, wat fijnstructuur toevoegt: het kenmerk
        bij \qty{15}{\micro m} is in werkelijkheid een \emph{band} van
        zo'n \qty{1}{\micro m} breed. Waarom doet de \emph{breedte} van de
        band ertoe voor een broeikasgas (denk aan wat er gebeurt zodra het
        midden van de band ondoorzichtig is)?
  \item Een trillingsaangeslagen CO$_2$ in lucht verliest zijn kwantum
        veel eerder door een botsing dan door straling (stralingslevensduur
        $\sim\qty{1}{s}$, botsingstijd $\sim\qty{e-9}{s}$): waar gaat de
        geabsorbeerde stralingsenergie werkelijk heen?
  \item Omgekeerd houdt lucht van \qty{288}{K} een thermische bezetting
        in $n = 1$ (vraag 8): wat doet die bezetting dat voor de
        energiebalans van belang is?
\end{enumerate}

\medskip\noindent\textbf{Deel III --- De stralingsboekhouding van de planeet.}
\begin{enumerate}[resume]
  \item De zon straalt als een lichaam van \qty{5800}{K}: bereken haar
        wienpiek. Onderschept de buiging van CO$_2$ veel zonlicht?
  \item De grond straalt als een lichaam van \qty{288}{K}: bereken de
        wienpiek en plaats \qty{15}{\micro m} op die thermische kromme.
  \item Leg het broeikasmechanisme in vier zinnen uit: zonlicht erin,
        thermisch infrarood eruit, onderschepping bij \qty{15}{\micro m},
        heruitzending zowel omhoog \emph{als} omlaag.
  \item De heruitzending die naar de ruimte ontsnapt komt uit hoge, koude
        lagen ($\sim\qty{220}{K}$): waarom verkleint uitzenden vanuit een
        koudere laag het uitgaande vermogen van de planeet bij die
        golflengten (denk eraan dat de thermische emissie met $T$ groeit)?
  \item Meer CO$_2$ duwt de uitzendende laag hoger en kouder: zeg in één
        zin waarom het oppervlak dan moet opwarmen om de boeken weer
        kloppend te maken.
  \item Het midden van de band is al ondoorzichtig; het effect van extra
        CO$_2$ werkt in de \emph{flanken} van de band, wat een
        logaritmische groei van de forcering met de concentratie geeft:
        breng dit in verband met vraag 11.
\end{enumerate}

\medskip\noindent\textbf{Deel IV --- Isotopen: de
kwantumvingerafdruk.}
\begin{enumerate}[resume]
  \item De frequentie van een oscillator schaalt als $\sqrt{k/\mu}$: bij
        de asymmetrische rek verandert het vervangen van $^{12}$C door
        $^{13}$C de effectieve massa; de waargenomen lijn verschuift van
        \qty{2349}{cm^{-1}} naar ongeveer \qty{2283}{cm^{-1}}. Ga de orde
        van grootte van deze verschuiving van $\approx 3\%$ na uit de
        massa's.
  \item Lasers die op deze twee lijnen zijn afgestemd tellen
        $^{13}$CO$_2$ en $^{12}$CO$_2$ afzonderlijk in een gasmonster:
        leg uit waarom de gekwantiseerde ladder zulke isotoopgescheiden
        detectie überhaupt mogelijk maakt.
  \item Planten geven de voorkeur aan de lichtste isotoop, dus is fossiele
        koolstof arm aan $^{13}$C: wat heeft de gemeten isotopenverhouding
        van CO$_2$ in de atmosfeer gedaan terwijl de concentratie steeg
        --- en wat bewijst dat over de \emph{bron} van het toegevoegde
        gas?
  \item Dezelfde natuurkunde bij \qty{15}{\micro m} werkt op Venus
        ($96\%$ CO$_2$, \qty{90}{bar}): wat illustreert haar oppervlak van
        \qty{737}{K}?
  \item Ook Mars ademt vrijwel zuivere CO$_2$, maar bij \qty{6}{mbar}, en
        is ijzig: wat toont het drietal Venus--aarde--Mars aan over welke
        grootheid de sterkte van het effect bepaalt?
  \item Vat het genoemde resultaat samen: een kwantum van \qty{83}{meV}
        --- de $\hbar\omega$ van een buigend drieatomig molecuul ---
        geparkeerd bij \qty{15}{\micro m} op het thermische spectrum van
        een planeet van \qty{288}{K}, geabsorbeerd, in nanoseconden
        gethermaliseerd en vanaf koude hoogten heruitgezonden, is het
        mechanisme waarmee $0.04\%$ van de lucht de temperatuur van de
        aarde bepaalt.
\end{enumerate}
\end{problem}
